% Generated draft: requires mathematical and source review.
\documentclass[11pt]{article}
\usepackage[margin=25mm]{geometry}
\usepackage{fontspec}
\setmainfont{Latin Modern Roman}
\usepackage{amsmath,amssymb,mathtools,unicode-math}
\setmathfont{Latin Modern Math}
\usepackage{enumitem,needspace,xurl,hyperref}
\setlength{\parindent}{0pt}
\setlength{\parskip}{5pt}
\setlength{\emergencystretch}{3em}
\newcommand{\N}{\mathbb{N}}
\newcommand{\Z}{\mathbb{Z}}
\newcommand{\Q}{\mathbb{Q}}
\newcommand{\R}{\mathbb{R}}
\newcommand{\C}{\mathbb{C}}
\newcommand{\F}{\mathbb{F}}
\newcommand{\A}{\mathbb{A}}
\newcommand{\PP}{\mathbb{P}}
\newcommand{\E}{\mathbb{E}}
\newcommand{\eps}{\varepsilon}
\newcommand{\dd}{\,\mathrm{d}}
\begin{document}
\section*{UnsolvedMath ID: 20000630}
\textbf{Source problem code:} catalogue entry
\subsection{Definitions and mathematical statement}
A Gaussian entire function (GAF) is interpreted here as a centered circular complex Gaussian process with almost surely entire sample functions. Let $N(A)$ count zeros in a bounded Borel set. A weak negative-correlation target is $\operatorname{Cov}(N(A),N(B))\le0$ for disjoint $A,B$; a pair-intensity formulation is $\rho_2(z,w)\le\rho_1(z)\rho_1(w)$ away from the diagonal. Strong negative association would be a further requirement. The original does not fix which meaning is intended.

\textbf{Original Question 5 (Balint Virag).} Find a Gaussian entire function with negatively correlated zeros. The source says there is such an example on the unit disk and relates the question to repulsion of random-polynomial zeros. It does not explicitly exclude finite-rank or finite-zero examples.

\textbf{Attached unreviewed literal example and obstruction.} For independent standard complex Gaussians $\xi_0,\xi_1$, the aid notes
\[G(z)=e^z(\xi_0+\xi_1z)\]
has exactly one zero almost surely. Its zero process is a rank-one projection determinantal process, and for disjoint sets the count covariance is minus the product of the two single-zero probabilities. The aid explicitly sets this finite-rank example aside as degenerate relative to the intended infinite-zero problem.

For a nondegenerate entire GAF with constant intensity $\rho_1=L/\pi$, $L>0$, the aid asserts that its zero law equals that of the planar GAF with covariance $e^{Lz\overline w}$, up to multiplication by a zero-free entire factor. This uses the stated centered circular GAF convention and does not impose only stationarity without a Gaussian-kernel hypothesis. Its normalized pair correlation is reported strictly greater than one when $t=L|z-w|^2/2\ge2$.

More quantitatively, for
\[A_L=D\left(0,\frac1{4\sqrt L}\right),\qquad B_L=D\left(\frac3{\sqrt L},\frac1{4\sqrt L}\right),\]
the aid claims
\[\operatorname{Cov}(N(A_L),N(B_L))\ge\frac1{128(e^{49/4}-1)}>0.\]
Thus the constant-positive-intensity class cannot satisfy even the weak disjoint-count negative-correlation target. This does not rule out inhomogeneous infinite-zero entire GAFs.

The aid further discusses the standard hyperbolic disk family. In the usual expanding-disk normalization its covariance is $(1-z\overline w/R^2)^{-\alpha_R}$. Keeping the determinantal parameter $\alpha_R=1$ while $R\to\infty$ loses local zero intensity; maintaining positive limiting intensity $\alpha_R/R^2\to L>0$ instead gives the planar kernel $e^{Lz\overline w}$ and the reported obstruction. This is about the scaling limit, not a claim that the parameter-one disk GAF itself has no zeros. The aid leaves the inhomogeneous infinite-zero question open. All kernel, correlation and status assertions require independent review.
\subsection{Short English statement}
The literal question permits a degenerate one-zero Gaussian example. For the intended infinite-zero problem, the aid rules out constant positive intensity through an explicit positive zero-count covariance. A disk-to-plane scaling discussion does not produce a negative-correlated entire example; the inhomogeneous case remains unresolved in the report.
\subsection{Sources}
\begin{itemize}
\item Local UnsolvedMath registry, permanent ID 20000630, source code catalogue entry; \texttt{aim-analysis-notes.json}, notes[161], tag \texttt{question}.
\item AIM workshop \emph{Random analytic functions}, Question 5 (Balint Virag), compiled by Swaminathan Sethuraman, April 17, 2006: \url{https://aimath.org/WWN/randomzeros/arcc1.pdf}. Record status \texttt{reconstructed\_unverified}, no clean statement. No PDF was consulted.
\item Attached machine-generated research attempt 1: result \texttt{partial\_result}, problem status at run \texttt{open}, rigor label \texttt{proved}, novelty confidence low. It is not peer reviewed and requires independent verification before citation. No authoritative bibliography for the quantitative obstruction or disk example is supplied.
\end{itemize}
\end{document}
